%%%%%%%%%%  scan idx 184  |  folio 171  %%%%%%%%%%
cohérent
définissant $S-U$. D'après (iii$'$), il existe $(i,n)\in I\times\mathbb N^+$ et une section $f$ de $F\otimes L_i^{\otimes n}$ au-dessus de $S$ telle que $f(x)\ne0$. On a $x\in S_f\subset U$, ce qui prouve (i).

\textbf{Définition 2.2.4.} Soit $S$ un schéma quasi-compact et quasi-séparé, et soit $(L_i)_{i\in I}$ une famille de faisceaux inversibles sur $S$. On dit que les $L_i$ forment une \emph{famille ample}\footnote{Cette notion, sous la forme (i$'$), a été introduite par M.~BORELLI et S.~L.~KLEIMAN ; la terminologie «~famille ample~» est proposée par le rédacteur.} si les conditions équivalentes de (2.2.3) sont vérifiées. (En particulier une famille réduite à un seul faisceau $L$ est ample si et seulement si $L$ est ample au sens usuel (EGA II 4.5.3)).

\textbf{Définition 2.2.5.} On dit qu'un schéma $S$ est \emph{divisoriel} si $S$ est quasi-compact et quasi-séparé et s'il existe une famille ample (2.2.4) de $\mathcal O_S$-Modules inversibles.

\textbf{Lemme 2.2.6.} Soit $X$ un schéma localement noethérien. Soit $U$ un ouvert de $X$ tel que l'inclusion $U\to X$ soit affine (par exemple $U$ affine et $X$ séparé). Alors $X-U$ est purement de codimension $0$ ou $1$ (i.e.\ en tout point maximal $x$ de $X-U$ on a $\dim(\mathcal O_{X,x})=0$ ou $1$), donc purement de codimension $1$ si $U$ contient les points maximaux de $X$.

\emph{Preuve.} Posons $Y=X-U$. Soit $x$ un point maximal de $Y$, faisons le changement de base $X'=\operatorname{Spec}(\mathcal O_{X,x})\to X$ : on obtient un diagramme cartésien :
\[
\begin{tikzcd}
U' \arrow[r] \arrow[d] & X' \arrow[d] & Y' \arrow[l] \arrow[d] \\
U \arrow[r] & X & Y \arrow[l]
\end{tikzcd},
\]
où $Y'=X'-U'$. Comme $x$ est un point maximal de $Y$, $Y'$ n'est autre que le point fermé de $X'$. D'autre part, comme $X'$ est affine, $U'$ est un ouvert affine de $X'$. Montrons que la dimension $n$ de $X'$ est égale à $0$ ou $1$. En effet, on sait (SGA 2 V 5.7) que $H^n_{Y'}(X',\mathcal O_{X'})\ne0$. Mais on a, si $n\ge2$,
\[
H^n_{Y'}(X',\mathcal O_{X'})\simeq H^{n-1}(U',\mathcal O_{U'}),
\]
or
\[
H^{n-1}(U',\mathcal O_{U'})=0,
\]
